\documentclass[12pt]{article}
\usepackage{../../ficha}

\begin{document}

\ficha{Geometria Analítica no Plano}[Distâncias, Mediatriz, Circunferência, Elipse, Retas e Regiões do Plano]

\instrucoes{Considera um referencial ortonormado $xOy$. Apresenta os cálculos e simplifica os radicais.}

\secao{Distância e Ponto Médio}{Calcula a distância entre os pontos e as coordenadas do ponto médio do segmento.}
\begin{exercicios}
  \item $A(1, 2)$ e $B(4, 6)$
  \item $A(-2, 3)$ e $B(4, -5)$
  \item $A(0, 0)$ e $B(5, 12)$
  \item $A(-1, -1)$ e $B(3, 1)$
  \item $A(2, -3)$ e $B(-4, 5)$
  \item Mostra que o triângulo de vértices $A(0,0)$, $B(4,0)$ e $C(2,3)$ é isósceles.
  \item Mostra que o triângulo de vértices $A(1,1)$, $B(4,1)$ e $C(4,5)$ é retângulo e calcula a sua área.
  \item O ponto médio de $[AB]$ é $M(3, 1)$ e $A(1, -2)$. Determina $B$.
\end{exercicios}

\secao{Mediatriz e Circunferência}{Determina a equação pedida.}
\begin{exercicios}
  \item A mediatriz do segmento de extremos $A(1, 1)$ e $B(5, 3)$.
  \item A mediatriz do segmento de extremos $A(-2, 0)$ e $B(4, 0)$.
  \item A circunferência de centro $(2, -1)$ e raio 3.
  \item A circunferência de centro na origem que passa por $(3, 4)$.
  \item A circunferência de diâmetro $[AB]$, com $A(-1, 2)$ e $B(5, 2)$.
  \item Indica o centro e o raio de $(x - 4)^2 + (y + 2)^2 = 25$.
  \item Indica o centro e o raio de $x^2 + y^2 - 6x + 4y - 3 = 0$.
  \item Verifica se o ponto $(1, 1)$ é interior, exterior ou pertence à circunferência $(x - 2)^2 + (y + 1)^2 = 4$.
\end{exercicios}

\secao{Elipse}{Resolve.}
\begin{exercicios}
  \item Indica os vértices, os focos e a excentricidade da elipse $\dfrac{x^2}{25} + \dfrac{y^2}{9} = 1$.
  \item Indica os vértices e os focos da elipse $\dfrac{x^2}{16} + \dfrac{y^2}{4} = 1$.
  \item Escreve a equação reduzida da elipse com focos $(\pm 4, 0)$ e eixo maior de comprimento 10.
  \item Escreve a equação reduzida da elipse com vértices $(\pm 6, 0)$ e $(0, \pm 2)$.
  \item O ponto $(3, y)$, com $y > 0$, pertence à elipse $\dfrac{x^2}{36} + \dfrac{y^2}{16} = 1$. Determina $y$.
  \item Define elipse como lugar geométrico dos pontos do plano.
\end{exercicios}

\secao{Retas no Plano}{Determina a equação reduzida da reta.}
\begin{exercicios}
  \item Passa em $A(1, 3)$ e tem declive 2.
  \item Passa em $A(2, 5)$ e $B(4, 9)$.
  \item Passa em $A(-1, 4)$ e $B(3, -4)$.
  \item Passa em $A(3, -2)$ e é paralela a $y = -\dfrac{1}{3}x + 1$.
  \item Passa em $A(0, 2)$ e tem vetor diretor $\vec u = (2, 3)$.
  \item É vertical e passa em $A(5, -1)$.
  \item Determina o declive e a inclinação (em graus) da reta $y = x - 4$.
  \item Determina a inclinação da reta $y = \sqrt{3}\,x + 2$.
  \item Determina o ponto de interseção das retas $y = 2x - 1$ e $y = -x + 5$.
\end{exercicios}

\secao{Regiões do Plano e Áreas}{Resolve.}
\begin{exercicios}
  \item Representa a região definida por $x \ge 0 \wedge y \ge 0 \wedge x + y \le 4$ e calcula a sua área.
  \item Representa a região definida por $x^2 + y^2 \le 9 \wedge y \ge 0$ e calcula a sua área.
  \item Escreve uma condição que defina o círculo de centro $(1, 2)$ e raio 2.
  \item Escreve uma condição que defina a coroa circular de centro na origem e raios 2 e 5.
  \item Calcula a área do triângulo de vértices $A(0, 0)$, $B(6, 0)$ e $C(2, 5)$.
  \item Calcula a área do triângulo limitado pelos eixos coordenados e pela reta $3x + 2y = 12$.
  \item Calcula a área do quadrilátero de vértices $(0, 0)$, $(4, 0)$, $(5, 3)$ e $(1, 3)$. Que quadrilátero é?
\end{exercicios}

\secao{Desafios}{Resolve os problemas.}
\begin{exercicios}
  \item Determina os pontos de interseção da circunferência $x^2 + y^2 = 25$ com a reta $y = x + 1$.
  \item Determina a equação da circunferência que passa nos pontos $(0, 0)$, $(6, 0)$ e $(0, 8)$.
  \item Determina o ponto do eixo $Ox$ equidistante de $A(1, 3)$ e $B(5, 1)$.
  \item Determina $k$ de modo que a reta $y = 2x + k$ passe pelo centro da circunferência $(x - 3)^2 + (y - 1)^2 = 4$.
  \item A órbita de um satélite é uma elipse com semieixos 10\,000\,km e 8\,000\,km, e a Terra num dos focos. Qual é a distância do centro da elipse ao foco?
\end{exercicios}

\end{document}
